
que nous noterons
@@M0@@

\textbf{Exemple 4.2.} Supposons que le schéma @@M1@@ soit de dimension @@M2@@. Prenons un élément @@M3@@, par exemple @@M4@@, où @@M5@@ est un sous-schéma intègre de dimension @@M6@@ de @@M7@@. Prenons aussi un polynôme homogène de poids @@M8@@
@@M9@@
où @@M10@@ est muni du poids @@M11@@ @@M12@@, par exemple un monôme @@M13@@ avec @@M14@@. Lorsqu'on substitue dans @@M15@@ pour chaque variable @@M16@@ @@M17@@ la classe de Chern
@@M18@@
on obtient un élément
@@M19@@
On appelle \emph{nombre de Chern de} @@M20@@ \emph{relativement à} @@M21@@ l'entier
@@M22@@
Lorsque @@M23@@, on omet la mention de @@M24@@ dans la terminologie précédente.

\subsection*{4.3.}
Sans hypothèse sur la dimension de @@M25@@, considérons une suite @@M26@@ de @@M27@@-modules inversibles et un @@M28@@-module cohérent @@M29@@ tel que @@M30@@. Posons @@M31@@ @@M32@@ et @@M33@@.

